\scanpage{27}
% transcription
entre donnée $\alpha$ et a), resp. $\beta$ et b) (ou c).
Dans le cas $C=$ ens. tot ordonnés $\ne\varnothing$, i.e. $\widehat C=$ ens. $\frac12$ simpl,
on a dans les $\operatorname{Hom}_{\cdot}(A,B)$ des relations d'homotopie.
Cela permet d'associer à $A$ des $n$-catégories pour tout $n$
(\unsure{Construction|construct} facile si les $\operatorname{Hom}_{\cdot}(A,B)$ sont des ens. $\frac12$ simpliciaux
de Kan \ldots)

\deleted{Hom} Définition de $\operatorname{Isom}_{\cdot}(A,B)$ et $\operatorname{Aut}_{\cdot}(A)$ (le premier
un groupe de $\widehat C$, le deuxième un pseudo-torseur sous
celui-ci). Notion d'homotopie de morphismes, d'isotopie
d'isomorphismes. Foncteurs \deleted{\unsure{de}} $\frac12$ simpliciaux de
catégories \inserted{à structure} semi-simpliciales \deleted{\unsure{Grp}} (Groupoïdes à structure
semi-simpliciale etc \ldots). On peut dire quand un
tel foncteur est une [$n$-]équivalence d'homotopie (
s'il induit \unsure{sur} les $\operatorname{Hom}_{\cdot}$ des ($n$)équivalences d'homotopie)
Sauf erreur, cela signifie aussi que le $n$-foncteur
entre les $n$-catégories associées est une $n$-équivalence.
\inserted{Cas des complexes d'une cat. abélienne. Lien avec ce qui précède via
Dold-Puppe.}

$2'$. \unsure{Tapis|Topis|Topos} de Quillen : interprétation des $n$-catégories
strictes comme structures $n$\inserted{fois}$\frac12$ simpliciales particulières.

3. $n$-catégories fibrées sur une catégorie, $n$-champs sur
un site un topos. \inserted{Ils forment une $(n+1)$-catégorie} Il faudrait voir

a) il y a une \inserted{$(n+1)$-}équivalence entre $n$-champs sur un
site, et $n$-champs sur le topos associé.

b) Toute $n$-catégorie fibrée sur un site admet un
$n$-champ enveloppant.

c) Existence d'images inverses \inserted{de $n$-champs} par morphismes
(de sites ou) de topos [ \uline{NB} La notion d'image directe
``est claire'' ; celle d'image inverse est censée s'en
déduire par \deleted{\unsure{une}} \inserted{\unsure{formule|formules} d'}adjonction
\[
\underline{\operatorname{Hom}}(A,f_*(B))\simeq\underline{\operatorname{Hom}}(f^*(A),B)
\]
[\unsure{$n$-équivalence|équivalence} \deleted{de} $n$-catégories \ldots].
% archive sheet number: 26
